\chapter{Credit default swaps}\label{ch:m2:credit-default-swaps}

Lehman Brothers a déposé le bilan le 15 septembre 2008. Des \omterm{def:m2:credit-default-swaps:cds}{credit default swaps}
portant sur Lehman, pour un notionnel brut d'environ 72~milliards de dollars, ont été enregistrés, et seulement environ
un tiers de cette protection a pu être acheté pour couvrir des obligations Lehman~; le reste était
détenu par des investisseurs qui ne possédaient aucune obligation à livrer. Le 10 octobre, quatorze dealers
ont organisé une enchère à deux étapes pour le compte de 358 parties~; elle a fixé le prix des obligations senior de Lehman à
8.625 pour 100, et chaque \omterm{def:m2:credit-default-swaps:cds}{vendeur de protection} ayant souscrit a payé 91.375 pour 100 du
notionnel, en espèces. Des positions compensatoires ont réduit les sommes échangées à environ
5.2~milliards de dollars. Ce chapitre explique le contrat à l'origine de ces paiements~: sa
forme standard depuis 2009, la manière dont il est coté et évalué à partir d'une \omterm{def:m2:credit-default-swaps:hazard}{intensité de défaut}, qui
décide qu'un défaut a eu lieu, comment l'enchère fixe le prix, et pourquoi l'obligation la moins chère
détermine le montant versé.

\section{Le contrat standard}

\begin{definition}[Credit default swap, acheteur et vendeur de protection, entité de référence]\label{def:m2:credit-default-swaps:cds}
Un \emph{credit default swap}\index{credit default swap} (CDS) est un contrat bilatéral dans lequel
l'\emph{acheteur de protection}\index{acheteur de protection} (protection buyer) verse une prime régulière sur un
montant notionnel jusqu'à l'échéance ou jusqu'à un \omterm{def:m2:credit-default-swaps:event}{événement de crédit} touchant un emprunteur désigné,
l'\emph{entité de référence}\index{entité de référence} (reference entity), selon ce qui survient en premier~; après un
\omterm{def:m2:credit-default-swaps:event}{événement de crédit}, le \emph{vendeur de protection}\index{vendeur de protection} (protection seller) verse à l'acheteur le
notionnel multiplié par un moins le \omterm{def:m2:corporate-bonds:covenant}{taux de recouvrement} de la dette de l'entité, fixé par le règlement-livraison.
\end{definition}

Un CDS est une assurance sur la dette d'un emprunteur, avec deux différences qui comptent~: l'acheteur
n'a pas besoin de détenir la dette, et le contrat peut se négocier et se compenser comme
n'importe quel swap. L'acheteur est short sur le crédit et le vendeur long, si bien qu'un \omterm{def:m2:credit-default-swaps:cds}{vendeur de
protection} a la même exposition au défaut qu'un détenteur d'obligations, sans avoir à financer l'obligation. Pour une protection de
10~millions de dollars à 100 points de base par an, l'acheteur verse 25\,000~dollars par trimestre~; si la
dette vaut 40 pour 100 après le défaut, le vendeur verse 6~millions de dollars
(voir la \cref{fig:m2:credit-default-swaps:flows}).

\begin{omfigure}
\centering
\begin{tikzpicture}[font=\small,
  box/.style={draw,thick,rounded corners=2pt,align=center,minimum height=1.1cm,minimum width=3.2cm,inner sep=3pt},
  arr/.style={-{Stealth[length=2.5mm]},thick}]
\node[box,draw=omDef,fill=omDef!8] (b) at (0,0) {acheteur de protection\\(short sur le crédit)};
\node[box,draw=omThm,fill=omThm!8] (s) at (9.4,0) {vendeur de protection\\(long sur le crédit)};
\node[box,draw=gray,fill=gray!8] (r) at (4.7,-3.2) {entité de référence\\(ses obligations et prêts)};
\draw[arr] ([yshift=3mm]b.east) -- node[above=2pt,font=\footnotesize,align=center] {prime~: $c \times$ notionnel par an,\\payée chaque trimestre jusqu'au défaut ou à l'échéance} ([yshift=3mm]s.west);
\draw[arr] ([yshift=-3mm]s.west) -- node[below=2pt,font=\footnotesize,align=center] {après un événement de crédit~:\\$(100 - \text{prix final})\%$ du notionnel} ([yshift=-3mm]b.east);
\draw[dashed,thick,gray] (r.north) -- (4.7,-1.45);
\node[font=\footnotesize,text=gray,align=right,anchor=east] at ([xshift=-4pt]r.west) {événement de crédit\\décidé par le comité de détermination};
\node[font=\footnotesize,text=gray,align=left,anchor=west] at ([xshift=4pt]r.east) {l'enchère fixe\\le prix final};
\end{tikzpicture}

\omcaption{Flux de trésorerie d'un \omterm{def:m2:credit-default-swaps:cds}{credit default swap}. L'acheteur verse une prime périodique
jusqu'à un \omterm{def:m2:credit-default-swaps:event}{événement de crédit} ou à l'échéance~; après un \omterm{def:m2:credit-default-swaps:event}{événement de crédit}, constaté par un comité, le
vendeur verse le notionnel multiplié par un moins le prix de la dette de l'entité fixé par une
enchère. Schéma de principe.}\label{fig:m2:credit-default-swaps:flows}
\end{omfigure}

\begin{definition}[Coupon standard, paiement initial (upfront)]\label{def:m2:credit-default-swaps:std}
Un \emph{coupon standard}\index{coupon standard} (standard coupon) est une prime périodique fixe, identique
pour tous les contrats d'une même catégorie, que le marché échange au lieu d'un coupon négocié
pour chaque transaction. Le \emph{paiement initial}\index{paiement initial} (upfront payment) est le montant
exprimé en pourcentage du notionnel, qu'une partie verse à l'autre au début de la transaction afin que le
contrat au coupon standard soit équitable au spread du marché~; l'\omterm{def:m2:credit-default-swaps:cds}{acheteur de protection}
le verse lorsque le spread est supérieur au coupon et le reçoit lorsqu'il est inférieur.
\end{definition}

Jusqu'en 2009, chaque transaction portait son propre coupon, le spread au pair du jour, de sorte que deux
transactions symétriques conclues à une semaine d'écart laissaient deux flux de primes différents qui
ne pouvaient être ni nettés ni compensés. À partir du 8 avril 2009, les contrats portant sur des entreprises nord-américaines
sont passés à des conditions standard~: un coupon fixe de 100 points de base pour les noms
investment grade, et de 500 pour les high yield, la différence étant versée en \omterm{def:m2:credit-default-swaps:std}{paiement initial}~;
échéances fixées aux dates de roll trimestrielles, les 20 mars, juin, septembre et décembre~;
les primes courant depuis la dernière date de roll précédant la transaction, le stub étant réglé avec le
\omterm{def:m2:credit-default-swaps:std}{paiement initial}~; la restructuration ne constitue plus un \omterm{def:m2:credit-default-swaps:event}{événement de crédit} pour ces contrats. Ce même changement a créé
les \omterm{def:m2:credit-default-swaps:event}{comités de détermination} et a rendu systématiques les enchères de la \cref{sec:m2:credit-default-swaps:events}.
Les contrats portant sur un même nom et une même échéance sont devenus identiques, et ont donc pu être nettés
et, plus tard, compensés par une contrepartie centrale.

\begin{dated}{2026-09}{La taille du marché des CDS}\label{dat:m2:credit-default-swaps:size}
À la fin de 2025, selon les statistiques de la BIS telles que résumées par l'ISDA~: notionnel en cours
de 11.0~billions de dollars, en hausse de 21.7~\% en un an, dont 4.6~billions de dollars en single-name et
6.4~billions de dollars sur plusieurs noms (indices et tranches, chapitre suivant)~; valeur de marché brute
de 273.7~milliards de dollars~; 70.2~\% du notionnel, soit 7.7~billions de dollars, compensés par
des contreparties centrales.
\end{dated}

\section{Paiement initial et prime périodique}

Pour convertir un spread en \omterm{def:m2:credit-default-swaps:std}{paiement initial}, et pour évaluer un contrat tout court, il
faut un modèle du moment où l'entité fait défaut.

\begin{definition}[Intensité de défaut, probabilité de survie]\label{def:m2:credit-default-swaps:hazard}
L'\emph{intensité de défaut}\index{intensité de défaut} (hazard rate) $\lambda(t)$ d'une \omterm{def:m2:credit-default-swaps:cds}{entité de référence} est son
taux instantané de défaut à l'instant $t$, sachant qu'elle a survécu jusqu'à $t$~: la
probabilité de défaut dans $[t, t + dt]$, sachant qu'il n'y a pas eu de défaut avant $t$, vaut $\lambda(t)\,dt$. La
\emph{probabilité de survie}\index{probabilité de survie} (survival probability) jusqu'à $T$ est $Q(T) =
\exp\bigl(-\int_0^T \lambda(u)\,du\bigr)$, qui vaut $e^{-\lambda T}$ lorsque l'intensité
de défaut est constante.
\end{definition}

Un CDS comporte deux jambes. La jambe de primes verse $c$ par an, trimestriellement, tant que
l'entité survit, plus la prime courue jusqu'au défaut~; sa valeur par unité de coupon est
l'\emph{annuité} risquée $A$ (celle du chapitre~9, chaque paiement étant pondéré par la
\omterm{def:m2:credit-default-swaps:hazard}{probabilité de survie} jusqu'à celui-ci). La jambe de protection verse $1 - R$ en cas de défaut, où $R$
est le \omterm{def:m2:corporate-bonds:covenant}{taux de recouvrement} du chapitre~21. Avec un taux $r$ constant et une \omterm{def:m2:credit-default-swaps:hazard}{intensité de défaut} $\lambda$ constante,
en supposant les défauts réglés à la fin du trimestre où ils surviennent, et en comptant un demi-
trimestre de prime courue,
\[
A = \sum_{k=1}^{4T} \tfrac14\, e^{-r t_k}\Bigl(Q(t_k) + \tfrac12\bigl(Q(t_{k-1}) - Q(t_k)\bigr)\Bigr),
\qquad
P = (1-R)\sum_{k=1}^{4T} e^{-r t_k}\bigl(Q(t_{k-1}) - Q(t_k)\bigr),
\]
avec $t_k = k/4$. Le spread au pair est $P/A$, la prime périodique qui rend le contrat
de valeur nulle.

\begin{proposition}[Le triangle de crédit]\label{prop:m2:credit-default-swaps:triangle}
Avec une \omterm{def:m2:credit-default-swaps:hazard}{intensité de défaut} $\lambda$ constante, un \omterm{def:m2:corporate-bonds:covenant}{taux de recouvrement} $R$ constant et des primes
versées en continu, le spread au pair vaut
\[
s = \lambda\,(1 - R),
\]
quels que soient le taux d'intérêt et l'échéance.
\end{proposition}

\begin{proof}
Avec des primes versées en continu, la jambe de primes vaut $s\int_0^T e^{-(r+\lambda)t}\,dt$
et la jambe de protection $(1-R)\int_0^T \lambda\,e^{-(r+\lambda)t}\,dt$, puisque le défaut
survient dans $[t, t+dt]$ avec la probabilité $\lambda e^{-\lambda t}dt$. Les intégrales
ne diffèrent que par le facteur constant $\lambda$, de sorte que les deux jambes sont égales lorsque $s =
\lambda(1-R)$.
\end{proof}

Avec des paiements trimestriels, le triangle n'est plus exact, mais il reste proche~: pour un spread de
100 points de base et un \omterm{def:m2:corporate-bonds:covenant}{taux de recouvrement} de 40~\%, l'\omterm{def:m2:credit-default-swaps:hazard}{intensité de défaut} qui valorise le
contrat à cinq ans est de 1.6667~\%, contre $0.01/0.6 = 1.6667\%$ pour le triangle, et
la \omterm{def:m2:credit-default-swaps:hazard}{probabilité de survie} à cinq ans est de 92.0~\%.

\begin{proposition}[Paiement initial à partir d'un spread coté]\label{prop:m2:credit-default-swaps:upfront}
Si un contrat de \omterm{def:m2:credit-default-swaps:std}{coupon standard} $c$ est coté à un spread $s$, et si $\lambda_s$ est
l'\omterm{def:m2:credit-default-swaps:hazard}{intensité de défaut} constante pour laquelle $s$ est le spread au pair, le \omterm{def:m2:credit-default-swaps:std}{paiement initial} (upfront) par unité
de notionnel, versé par l'\omterm{def:m2:credit-default-swaps:cds}{acheteur de protection}, est
\[
U = (s - c)\,A(\lambda_s).
\]
\end{proposition}

\begin{proof}
Pour $\lambda_s$, la jambe de protection vaut $s\,A(\lambda_s)$, par définition du
spread au pair. Le contrat de coupon $c$ vaut $P - cA = (s - c)A$ pour l'acheteur,
qui est le montant que l'acheteur verse d'avance pour rendre la transaction équitable.
\end{proof}

\begin{example}[Deux cotations, deux conventions]\label{ex:m2:credit-default-swaps:upfront}
Avec un taux constant de 4~\% et un \omterm{def:m2:corporate-bonds:covenant}{taux de recouvrement} de 40~\%, l'annuité risquée à cinq ans vaut 4.332
pour un spread de 100 points de base et 4.165 pour 200. Un émetteur investment grade coté à
200 points de base, sur le coupon de 100 points de base, coûte à l'acheteur $1\% \times 4.165 =
4.16\%$ de \omterm{def:m2:credit-default-swaps:std}{paiement initial} (upfront), soit 416\,000~USD sur un notionnel de 10~millions de dollars. Un émetteur high yield coté à 300
points de base sur le coupon de 500~: l'acheteur touche 8.01~\% en \omterm{def:m2:credit-default-swaps:std}{paiement initial}~; à 800, il verse
9.98~\%, et à 1\,200, 20.28~\%. Les émetteurs investment grade se cotent en spread,
les émetteurs high yield en points de \omterm{def:m2:credit-default-swaps:std}{paiement initial}.
\end{example}

La \cref{fig:m2:credit-default-swaps:upfront} montre le \omterm{def:m2:credit-default-swaps:std}{paiement initial} en fonction du spread pour
les deux coupons~: elle croise zéro au niveau du coupon, et elle s'incurve parce que l'annuité
se réduit quand l'\omterm{def:m2:credit-default-swaps:hazard}{intensité de défaut} augmente. La pente au voisinage d'un spread est l'annuité, donc 10~millions de dollars de notionnel
de protection à cinq ans autour de 100 points de base rapportent environ 4\,300~USD pour chaque
point de base dont le spread s'élargit. La \cref{fig:m2:credit-default-swaps:survival} montre ce que
les spreads signifient pour la survie.

\begin{omfigure}
\begin{tikzpicture}
\begin{axis}[width=11cm,height=5cm,
  xlabel={spread coté (pb)},ylabel={paiement initial versé par l'acheteur (\%)},
  tick label style={font=\small},label style={font=\small},
  xmin=0,xmax=1500,ymin=-25,ymax=40,grid=major,grid style={gray!25},
  legend columns=2,legend style={font=\footnotesize,at={(0.5,-0.3)},anchor=north,draw=none,fill=none,/tikz/every even column/.append style={column sep=8pt}},
  table/col sep=comma]
\addplot[omDef,very thick] table[x=spread,y=c100]{figdata/markets-2/23-credit-default-swaps/upfront.csv};
\addplot[omThm,very thick] table[x=spread,y=c500]{figdata/markets-2/23-credit-default-swaps/upfront.csv};
\addplot[gray,thin,forget plot] coordinates {(0,0) (1500,0)};
\addplot[only marks,mark=*,mark size=1.6pt,black,forget plot] coordinates {(200,4.16) (300,-8.01) (800,9.98) (1200,20.28)};
\legend{coupon de 100 pb,coupon de 500 pb}
\end{axis}
\end{tikzpicture}

\omcaption{\omterm{def:m2:credit-default-swaps:std}{Paiement initial} d'un contrat à cinq ans, en pourcentage du notionnel, en fonction
du spread coté, pour les deux \omterm{def:m2:credit-default-swaps:std}{coupons standard} (taux constant de 4~\%, \omterm{def:m2:corporate-bonds:covenant}{taux de recouvrement} de 40~\%). Le
\omterm{def:m2:credit-default-swaps:std}{paiement initial} est nul là où le spread égale le coupon, et négatif, versé à l'acheteur,
en dessous~; les points correspondent à l'\cref{ex:m2:credit-default-swaps:upfront}. Illustratif~; données~:
le tutoriel du chapitre.}\label{fig:m2:credit-default-swaps:upfront}
\end{omfigure}

\begin{omfigure}
\begin{tikzpicture}
\begin{axis}[width=11cm,height=4.8cm,
  xlabel={années},ylabel={probabilité de survie},
  tick label style={font=\small},label style={font=\small},
  xmin=0,xmax=10,ymin=0.2,ymax=1.02,grid=major,grid style={gray!25},
  yticklabel style={/pgf/number format/.cd,fixed,precision=1},
  legend columns=3,legend style={font=\footnotesize,at={(0.5,-0.3)},anchor=north,draw=none,fill=none,/tikz/every even column/.append style={column sep=8pt}},
  table/col sep=comma]
\addplot[omDef,very thick] table[x=t,y=s100]{figdata/markets-2/23-credit-default-swaps/survival.csv};
\addplot[omThm,very thick] table[x=t,y=s300]{figdata/markets-2/23-credit-default-swaps/survival.csv};
\addplot[omProp,very thick] table[x=t,y=s800]{figdata/markets-2/23-credit-default-swaps/survival.csv};
\legend{100 pb,300 pb,800 pb}
\end{axis}
\end{tikzpicture}

\omcaption{\omterm{def:m2:credit-default-swaps:hazard}{Probabilités de survie} implicites des spreads à cinq ans de 100, 300 et 800
points de base, avec une \omterm{def:m2:credit-default-swaps:hazard}{intensité de défaut} constante et un \omterm{def:m2:corporate-bonds:covenant}{taux de recouvrement} de 40~\%. Les \omterm{def:m2:credit-default-swaps:hazard}{intensités de défaut} sont
proches du spread divisé par $1 - R$~: environ 1.7~\%, 5~\% et 13~\% par an. Ce sont
des probabilités neutres au risque, qui incluent une prime pour supporter le risque de défaut.
Illustratif~; données~: le tutoriel du chapitre.}\label{fig:m2:credit-default-swaps:survival}
\end{omfigure}

Les probabilités sont celles qui valorisent le contrat, et non des prévisions. Le chapitre~21
a décomposé le spread d'une obligation en perte attendue et en prime de risque et d'illiquidité~;
la même décomposition vaut ici~: une \omterm{def:m2:credit-default-swaps:hazard}{intensité de défaut} déduite d'un spread est donc plus élevée que le
taux de défaut que laisse supposer l'historique.

\section{Événements de crédit et comités de détermination}\label{sec:m2:credit-default-swaps:events}

\begin{definition}[Événement de crédit, comité de détermination]\label{def:m2:credit-default-swaps:event}
Un \emph{événement de crédit}\index{evenement@événement de crédit} est l'un des événements listés dans un contrat CDS
dont la survenance déclenche le paiement de protection, comme la faillite de l'\omterm{def:m2:credit-default-swaps:cds}{entité de
référence} ou son défaut de paiement. Un \emph{comité de
détermination}\index{comité de détermination} (DC) est l'organe formé des dealers et des investisseurs
qui décide, pour tous les contrats standard à la fois, si et quand un événement de crédit
s'est produit et comment il est réglé.
\end{definition}

La documentation de l'ISDA énumère les \omterm{def:m2:credit-default-swaps:event}{événements de crédit} qu'une transaction peut préciser~:
faillite, défaut de paiement, restructuration, accélération d'obligation, défaut d'obligation
répudiation ou moratoire~; le contrat précise lesquels s'appliquent. La faillite et le défaut
de paiement sont les plus courants pour les entreprises, la restructuration concerne les banques et les États souverains,
et savoir si un événement donné remplit les conditions requises est souvent une question d'interprétation. Avant
2009, chaque paire de contreparties décidait pour son propre compte, et deux transactions de sens contraire pouvaient
être déclenchées différemment.

La réforme de 2009 a créé un comité pour chacune des cinq régions (les Amériques, l'Asie
hors Japon, le Japon, l'Australie et la Nouvelle-Zélande, et l'Europe, le Moyen-Orient et
l'Afrique), comptant à sa création dix dealers et cinq autres membres votants, et l'ISDA en tant que
secrétaire. Toute partie à une transaction couverte peut poser une question~; le comité décide
si un \omterm{def:m2:credit-default-swaps:event}{événement de crédit} s'est produit et à quelle date, s'il faut organiser une enchère, quelles
obligations peuvent être livrées, et si l'entité a un successeur. La plupart des décisions
requièrent 80~\% des membres, à défaut de quoi un collège externe tranche, et elles lient toutes
les transactions couvertes. Une période de rétrospection de 60~jours a remplacé la date de début propre à chaque transaction, de sorte que
des transactions de sens contraire conclues à des dates différentes soient déclenchées par les mêmes événements.

Les règles évoluent au gré des cas. En 2013, l'État néerlandais a nationalisé la banque SNS
et exproprié ses obligations subordonnées, un événement que les définitions de l'époque ne
couvraient pas clairement. Les définitions de 2014, en vigueur à partir du 22 septembre 2014, ont ajouté un
\emph{intervention gouvernementale}, nouvel \omterm{def:m2:credit-default-swaps:event}{événement de crédit} pour les entités financières hors des
États-Unis et permettent aux \omterm{def:m2:credit-default-swaps:cds}{acheteurs de protection} de livrer tous les actifs vers lesquels un
renflouement interne (bail-in) ou une restructuration convertit les obligations.

\section{Enchères de règlement-livraison}

Après un \omterm{def:m2:credit-default-swaps:event}{événement de crédit}, le paiement vaut $1 - R$ par unité de notionnel, et quelqu'un doit
fixer $R$. Le règlement-livraison physique, méthode d'origine, oblige l'acheteur à livrer des obligations de
l'entité et à recevoir le pair~; mais lorsque l'essentiel de la protection est détenu par des investisseurs sans obligations, comme
avec Lehman, ces acheteurs doivent acheter des obligations pour les livrer, ce qui fait monter leur cours et baisser le montant versé.
Le règlement en espèces suppose un prix unique pour des obligations qui se négocient à peine. L'enchère
le fournit.

\begin{definition}[Prix final d'enchère]\label{def:m2:credit-default-swaps:final}
Le \emph{prix final d'enchère}\index{prix final d'enchère} d'un \omterm{def:m2:credit-default-swaps:event}{événement de crédit} est le
prix pour 100 de pair de la dette livrable de l'entité, fixé par l'enchère de
règlement-livraison, auquel tout CDS couvert par l'enchère est réglé en espèces et toute
demande de règlement-livraison physique formulée lors de l'enchère est exécutée.
\end{definition}

L'enchère comporte deux étapes. Dans la première, chaque dealer participant soumet un cours acheteur
et un cours vendeur écartés de 2 au plus, pour une taille de cotation fixe d'obligations, ainsi que
des demandes de règlement-livraison physique, pour son propre compte et celui de ses clients, d'achat ou de vente d'obligations au
prix final, chacune n'excédant pas la position CDS de la partie et de même sens.
Les administrateurs retirent les paires croisées, le cours acheteur le plus élevé contre le cours vendeur le plus bas, jusqu'à ce qu'aucun
cours acheteur n'atteigne un cours vendeur, puis calculent la moyenne de la meilleure moitié des cours acheteurs et vendeurs restants,
arrondie au huitième le plus proche~: le \emph{point médian du marché interne}. Le solde des
demandes constitue les \emph{positions ouvertes}. Un dealer dont la cotation a croisé du mauvais côté
par rapport au point médian, compte tenu du sens des positions ouvertes, paie une pénalité.

Dans la seconde étape, les participants soumettent des ordres à cours limité du côté opposé aux
positions ouvertes~; les ordres sont exécutés au meilleur prix d'abord jusqu'à ce que les
positions ouvertes soient couvertes, et le dernier ordre exécuté fixe le prix final. Le prix ne peut pas être supérieur de 1 au
point médian lorsque les positions ouvertes sont à la vente, ou inférieur de 1 lorsqu'elles sont à l'achat.

\begin{omfigure}
\begin{tikzpicture}
\begin{axis}[width=11cm,height=5cm,
  xlabel={dealer (trié par cours acheteur)},ylabel={prix pour 100},
  tick label style={font=\small},label style={font=\small},
  xmin=0.3,xmax=14.7,ymin=7.5,ymax=12.5,grid=major,grid style={gray!25},
  xtick={1,2,3,4,5,6,7,8,9,10,11,12,13,14},
  legend columns=4,legend style={font=\footnotesize,at={(0.5,-0.3)},anchor=north,draw=none,fill=none,/tikz/every even column/.append style={column sep=8pt}},
  table/col sep=comma]
\addplot[only marks,mark=triangle*,mark size=2.6pt,omDef] table[x=dealer,y=bid]{figdata/markets-2/23-credit-default-swaps/lehman.csv};
\addplot[only marks,mark=triangle*,mark options={rotate=180},mark size=2.6pt,omThm] table[x=dealer,y=offer]{figdata/markets-2/23-credit-default-swaps/lehman.csv};
\addplot[omProp,very thick,dashed] coordinates {(0.3,9.75) (14.7,9.75)};
\addplot[black,very thick] coordinates {(0.3,8.625) (14.7,8.625)};
\legend{cours acheteur,cours vendeur,point médian 9.75,prix final 8.625}
\end{axis}
\end{tikzpicture}

\omcaption{Première étape de l'enchère Lehman Brothers du 10 octobre 2008~: cours acheteur et cours vendeur de chacun des
quatorze dealers pour 5~millions de dollars d'obligations. Une paire a croisé (un cours acheteur et
un cours vendeur à 10)~; la moyenne de la meilleure moitié des cours restants a donné un point
médian du marché interne de 9.75. Les positions ouvertes étaient de 4.92~milliards de dollars à la vente, et la seconde
étape a fixé le prix final à 8.625. Données~: Helwege, Maurer, Sarkar et Wang, FRBNY Staff
Report 372, encadré 1.}\label{fig:m2:credit-default-swaps:lehman}
\end{omfigure}

\begin{example}[Lehman Brothers]\label{ex:m2:credit-default-swaps:lehman}
À la première étape, le cours acheteur le plus élevé, 10, a rencontré un cours vendeur de 10 et les deux ont été retirés~; la
meilleure moitié, soit sept des treize cours acheteurs et vendeurs restants, a donné une moyenne de 9.804, arrondie à
9.75. Les demandes de vente d'obligations totalisaient 5.69~milliards de dollars et les demandes d'achat 0.77~milliard de dollars~:
des positions ouvertes de 4.92~milliards de dollars à la vente, plafonnant le prix final à 10.75. Le
dealer dont le cours acheteur à 10 a croisé, au-dessus du point médian alors que les positions ouvertes étaient à la
vente, a payé $5\,\text{millions} \times (10 - 9.75)\% = \text{USD}~12\,500$. Dans la seconde
étape, 453 ordres d'achat ont été soumis, de 10.75 à 0.125~; ceux à 8.625 ou
plus élevés ont totalisé les positions ouvertes, et le prix final était de 8.625
(la \cref{fig:m2:credit-default-swaps:lehman}).
\end{example}

L'enchère fait deux choses à la fois. Elle fixe un prix unique, si bien qu'un investisseur couvert
par un contrat sur nom unique contre un contrat indiciel ne se retrouve pas avec des \omterm{def:m2:corporate-bonds:covenant}{taux de
recouvrement} différents selon le nom. Elle transforme aussi le règlement-livraison physique en une transaction sur le solde net~: un
\omterm{def:m2:credit-default-swaps:cds}{acheteur de protection} qui détient des obligations et veut s'en défaire demande à les vendre au
prix final et finit au pair, tandis que tous les autres se règlent en espèces. Pour Lehman, les
72~milliards de dollars de protection brute sont devenus 5.2~milliards de dollars de paiements nets, car la plupart
des dealers avaient acheté et vendu de la protection sur ce nom.

Le prix n'est pas garanti de coller au marché obligataire. Les obligations Lehman se sont négociées à 34 le
jour de la faillite, 13 la veille de l'enchère et environ 10 le jour même~; le
prix final, tiré vers le bas par les importantes positions ouvertes à la vente, était un peu inférieur. Lors de l'enchère Fannie Mae de 2008, une forte demande dans la
première étape a poussé le prix final de la dette subordonnée au-dessus de celui de la dette senior.

\section{L'option du moins cher à livrer}

Un CDS ne protège pas une seule obligation mais une catégorie de la dette de l'entité~: toute obligation du
rang de séniorité, devise et forme requis est livrable. Après un défaut, les obligations de même
rang sont des créances sur la même masse de faillite et devraient se négocier à des prix voisins, mais
pas toujours~: une obligation longue à faible coupon et une obligation courte à coupon élevé peuvent
diverger avant que la faillite ne soit réglée, et lors d'une restructuration les obligations peuvent être
traitées différemment. L'\omterm{def:m2:credit-default-swaps:cds}{acheteur de protection} qui opte pour le règlement-livraison physique livre la moins chère
d'entre elles, et lors de l'enchère les vendeurs d'obligations livrent eux aussi la moins chère, si bien que le
prix final est fixé par la moins chère des obligations livrables~: l'option du \omterm{def:m2:bond-futures:ctd}{moins cher à livrer} du
chapitre~6, appliquée désormais au crédit. L'option vaut le plus lorsque les obligations livrables diffèrent de prix,
ce qui est précisément le cas où elle compte le plus~: les restructurations qui laissent certaines obligations à longue échéance et
à faible coupon, ou les renflouements internes qui convertissent les obligations en d'autres actifs.

\section{Tutoriel~: évaluer un CDS et rejouer une enchère}\label{tut:m2:credit-default-swaps:cds}

\begin{tutorial}
\textbf{Objectif.} Évaluer un CDS à \omterm{def:m2:credit-default-swaps:hazard}{intensité de défaut} constante, convertir les spreads cotés en paiements initiaux (upfront)
pour les deux \omterm{def:m2:credit-default-swaps:std}{coupons standard}, vérifier le triangle de crédit et rejouer
l'enchère de Lehman.
\textbf{État final~:} \cref{fig:m2:credit-default-swaps:upfront,fig:m2:credit-default-swaps:survival,fig:m2:credit-default-swaps:lehman,fig:m2:credit-default-swaps:orders},
\cref{ex:m2:credit-default-swaps:upfront,ex:m2:credit-default-swaps:lehman} et les
chiffres du problème du week-end.
\begin{tutsteps}
\item \textbf{Évaluation}~: les deux jambes, le spread au pair, le \omterm{def:m2:credit-default-swaps:std}{paiement initial} (upfront),
et l'\omterm{def:m2:credit-default-swaps:hazard}{intensité de défaut} implicite dans un spread.
\omcode{code/firm/cds/firm_cds.py}{14}{53}{Un CDS à intensité de défaut constante~: jambes, spread au pair, paiement initial (upfront) et intensité implicite.}\label{lst:m2:credit-default-swaps:price}
\item \textbf{L'enchère}~: point médian du marché intérieur (inside market midpoint), positions ouvertes, prix final et règlement
en espèces.
\omcode{code/firm/cds/firm_cds.py}{58}{96}{Enchère de règlement-livraison en deux étapes.}\label{lst:m2:credit-default-swaps:auction}
\item \textbf{Exécuter} \texttt{cds\_demo.tutorial()}, \texttt{cds\_demo.problem()} et
\texttt{fig\_cds.py}~; construire et tester les versions C++20 et Rust dans
\texttt{code/firm/cds/cpp} et \texttt{code/firm/cds/rust}.
\end{tutsteps}
\textbf{Pour aller plus loin.} Remplacer l'intensité constante par une courbe bootstrappée à partir des spreads à un,
trois, cinq et sept ans, et voir de combien varie le \omterm{def:m2:credit-default-swaps:std}{paiement initial} à cinq ans~;
puis faire varier le \omterm{def:m2:corporate-bonds:covenant}{taux de recouvrement} et vérifier que le \omterm{def:m2:credit-default-swaps:std}{paiement initial} bouge à peine lorsque l'\omterm{def:m2:credit-default-swaps:hazard}{intensité de défaut}
est recalibrée sur le même spread.
\end{tutorial}

\section{Construction~: le calculateur de prix des CDS}\label{bld:m2:credit-default-swaps:cds}

\begin{build}
\bfield{Rôle} La petite entreprise cote et valorise des CDS sur ses émetteurs obligataires~: paiements initiaux (upfront)
et conversions en primes courantes, \omterm{def:m2:credit-default-swaps:hazard}{intensités de défaut} pour ses modèles de crédit, et le calcul de
l'enchère lorsqu'un nom fait défaut. Comme le calculateur d'obligations du chapitre~3, il existe en
Python, C++20 et Rust, avec les mêmes tests.
\bfield{Interface} \texttt{Cds(maturity, coupon, recovery, freq)}~;
\texttt{legs(cds, lam, r)}~; \texttt{par\_spread}~; \texttt{upfront}~;
\texttt{hazard\_from\_spread(cds, spread, r)}~; \texttt{upfront\_from\_spread}~;
\texttt{inside\_market\_midpoint(bids, offers)}~; \texttt{open\_interest(requests)}~;
\texttt{final\_price(imm, oi, orders)}~; \texttt{cash\_settlement(notional, price)}.
\bfield{Règles} \omterm{def:m2:credit-default-swaps:hazard}{Intensité de défaut} constante et taux à capitalisation continue constant~; primes trimestrielles
avec une demi-période de prime courue en cas de défaut~; défauts payés à la fin du
trimestre~; prix d'enchère pour 100 de pair, point médian arrondi au huitième, prix de seconde étape
plafonné à un point au-delà du point médian~; les cotations
de première étape des dealers ne sont pas reprises en seconde étape.
\bfield{Tests d'acceptation} \texttt{code/firm/cds/tests/}, \texttt{cpp/firm\_cds\_test.cpp},
\texttt{rust/src/lib.rs}~: le spread au pair se retrouve après passage par l'intensité implicite~; le
\omterm{def:m2:credit-default-swaps:std}{paiement initial} est nul au coupon et change de signe de part et d'autre~; la première étape de Lehman donne
9.75 et des positions ouvertes de 4\,920~; le plafond est atteint à 10.75.
\bfield{Pour aller plus loin} Courbes d'intensité constantes par morceaux, bootstrappées à partir d'une structure par terme
de spreads~; dates exactes et intérêts courus sur le calendrier des dates de roll~; un CDS sur un indice
(chapitre suivant)~; le risque d'un livre de CDS par nom et par échéance (livre~6).
\end{build}

\begin{omsources}
\item J.~Helwege, S.~Maurer, A.~Sarkar et Y.~Wang, «~Credit default swap auctions~»,
Federal Reserve Bank of New York Staff Report 372, mai 2009.
\item Paul, Weiss, «~The Big Bang Protocol and a new structural framework for credit
default swaps~», mémorandum client, 24 mars 2009.
\item F.~Carruzzo, «~New ISDA 2014 Credit Derivatives Definitions~», Harvard Law School
Forum on Corporate Governance, août 2014.
\item ISDA, Disclosure Annex for Credit Derivative Transactions.
\item ISDA, «~Key trends in the size and composition of OTC derivatives markets in the
second half of 2025~», juillet 2026.
\end{omsources}

\section{Exercices}

\begin{exercise}[$\star$]\label{exo:m2:credit-default-swaps:1}
Un investisseur achète pour 10~millions de dollars de protection à 100 points de base par an, payés
trimestriellement. Que verse-t-il à chaque trimestre, et que reçoit-il si l'entité
fait défaut et que le \omterm{def:m2:credit-default-swaps:final}{prix final d'enchère} vaut 40~?
\end{exercise}

\begin{exercise}[$\star$]\label{exo:m2:credit-default-swaps:2}
Le spread à cinq ans d'un nom est de 300 points de base et le \omterm{def:m2:corporate-bonds:covenant}{taux de recouvrement} est pris égal à 40~\%. Utiliser
le triangle de crédit pour estimer son \omterm{def:m2:credit-default-swaps:hazard}{intensité de défaut} et sa \omterm{def:m2:credit-default-swaps:hazard}{probabilité de survie} à cinq ans.
\end{exercise}

\begin{exercise}[$\star$]\label{exo:m2:credit-default-swaps:3}
Énumérer quatre questions qu'un \omterm{def:m2:credit-default-swaps:event}{comité de détermination} tranche après un éventuel \omterm{def:m2:credit-default-swaps:event}{événement de crédit}.
\end{exercise}

\begin{exercise}[$\star\star$]\label{exo:m2:credit-default-swaps:4}
Un contrat à cinq ans sur un nom investment grade est coté à 50 points de base. Avec les
hypothèses du chapitre, quel est le \omterm{def:m2:credit-default-swaps:std}{paiement initial} sur 10~millions de dollars, et qui le verse~?
\end{exercise}

\begin{exercise}[$\star\star$]\label{exo:m2:credit-default-swaps:5}
Pourquoi les \omterm{def:m2:credit-default-swaps:std}{coupons standard} ont-ils rendu possibles la compensation multilatérale et la compensation par chambre des CDS, alors que les contrats
portant chacun son propre spread au pair ne le permettaient pas~?
\end{exercise}

\begin{exercise}[$\star\star$]\label{exo:m2:credit-default-swaps:6}
Après un défaut, deux obligations de même rang de séniorité sont livrables, l'une négociée à 30 et
l'autre à 35. Que reçoit un \omterm{def:m2:credit-default-swaps:cds}{acheteur de protection} de 10~millions de dollars par règlement-livraison physique
et quel prix l'enchère doit-elle refléter~?
\end{exercise}

\begin{exercise}[$\star\star\star$]\label{exo:m2:credit-default-swaps:7}
\emph{Programmation.} Avec \texttt{firm\_cds}, trouver l'\omterm{def:m2:credit-default-swaps:hazard}{intensité de défaut} implicite par les spreads à cinq ans
de 100, 300 et 800 points de base, et comparer chacune avec le triangle de crédit.
Expliquer l'écart.
\end{exercise}

\begin{exercise}[$\star\star\star$]\label{exo:m2:credit-default-swaps:8}
\emph{Trouver l'erreur.} «~Le spread d'un CDS est la probabilité de défaut du marché par
an~: un nom à 300 points de base a une probabilité de 3~\% de faire défaut chaque année.~» Corriger cette affirmation.
\end{exercise}

\section{Problème~: L'enchère}

\begin{problem}[{problème du week-end --- les deux étapes d'une enchère de règlement-livraison}]\label{pb:m2:credit-default-swaps:1}
Rejouer l'enchère de Lehman Brothers du 10 octobre 2008 avec les soumissions de première étape
de la \cref{fig:m2:credit-default-swaps:lehman}~: quatorze dealers ont soumis leurs cours acheteurs et vendeurs pour
5~millions de dollars~USD d'obligations, et des demandes de règlement-livraison physique (en millions de dollars~USD) visant à acheter 130, 612
et 30 et à vendre 755, 870, 141, 480, 464, 1\,470, 390, 574, 191, 170 et 187. Seuls
l'intervalle et la médiane des ordres de la seconde étape ont été publiés~; utiliser le carnet illustratif
d'ordres d'achat à cours limité de la \cref{fig:m2:credit-default-swaps:orders}, cohérent avec
le résultat publié. Quatre participants illustratifs règlent leurs positions~: A, un fonds de pension, détient
20~millions de dollars~USD d'obligations et 20~millions de dollars~USD de protection~; B, un hedge fund, a acheté 50
millions de dollars~USD de protection et ne détient aucune obligation~; C, un dealer, a vendu 100
millions de dollars~USD de protection et a acheté 60~millions de dollars~USD~; D, un assureur, a vendu 30~millions de dollars~USD de protection.

\begin{omfigure}
\begin{tikzpicture}
\begin{axis}[width=11cm,height=4.8cm,
  xlabel={cours acheteur limite (prix pour 100)},ylabel={taille cumulée (Md USD)},
  tick label style={font=\small},label style={font=\small},
  xmin=7.3,xmax=11,ymin=0,ymax=14,grid=major,grid style={gray!25},
  x dir=reverse,
  legend columns=2,legend style={font=\footnotesize,at={(0.5,-0.3)},anchor=north,draw=none,fill=none,/tikz/every even column/.append style={column sep=8pt}},
  table/col sep=comma]
\addplot[omDef,very thick,const plot,mark=*,mark size=1.4pt] table[x=price,y expr=\thisrow{cum}/1000]{figdata/markets-2/23-credit-default-swaps/orders.csv};
\addplot[omThm,very thick,dashed] coordinates {(11,4.92) (7.3,4.92)};
\addplot[black,thick,forget plot] coordinates {(8.625,0) (8.625,14)};
\node[font=\footnotesize,anchor=east] at (axis cs:8.66,11.5) {prix final 8.625};
\legend{cours acheteurs exécutés à partir du plus haut,intérêt ouvert 4.92}
\end{axis}
\end{tikzpicture}

\omcaption{La seconde étape de l'enchère sur un carnet illustratif d'ordres d'achat à cours limité
sur les obligations~: les cours acheteurs sont exécutés à partir du prix le plus élevé jusqu'à rencontrer l'intérêt ouvert
de 4.92~milliards de dollars~USD à vendre, ce qui se produit à 8.625, le prix final
publié. Les cours acheteurs inférieurs à 7.5 sortent du graphique. Carnet illustratif~; intérêt ouvert et résultat~:
FRBNY Staff Report 372.}\label{fig:m2:credit-default-swaps:orders}
\end{omfigure}

\medskip\noindent\textbf{Partie I --- La première étape.}
\begin{enumerate}
\item Quels cours acheteur et vendeur se croisent, et sont retirés~?
\item Calculer la moyenne de la meilleure moitié des cours acheteurs et vendeurs restants, et l'arrondir au
milieu interne du marché.
\item Calculer l'intérêt ouvert et son sens.
\item Quel dealer paie une pénalité, et de combien~?
\item Pourquoi l'intérêt ouvert portait-il sur la vente d'obligations~?
\end{enumerate}

\medskip\noindent\textbf{Partie II --- La seconde étape.}
\begin{enumerate}[resume]
\item Quels ordres sont soumis, et quel est le prix le plus élevé autorisé~?
\item Remplir le carnet à partir du cours acheteur le plus élevé et déterminer le prix final.
\item Quel serait le prix final si tous les ordres étaient à 12~?
\item Si seulement 1 milliard de dollars~USD avait été proposé à 8.625, quel serait le prix final~?
\item De combien le prix final s'est-il écarté du milieu, et pourquoi~?
\end{enumerate}

\medskip\noindent\textbf{Partie III --- Règlement-livraison.}
\begin{enumerate}[resume]
\item Un cash règle~: que reçoit-il, et que vaut son obligation au
prix final~?
\item A demande plutôt à vendre ses 20~millions de dollars~USD d'obligations dans l'enchère~: que
reçoit-il au total~?
\item Que reçoit B~?
\item Que paie C~?
\item D demande à acheter 30~millions de dollars~USD d'obligations dans l'enchère~: que paie-t-il, et que
possède-t-il ensuite~?
\end{enumerate}

\medskip\noindent\textbf{Partie IV --- Jugement.}
\begin{enumerate}[resume]
\item Pourquoi 72~milliards de dollars~USD de protection brute se sont-ils réduits à seulement
5.2~milliards de dollars~USD de paiements~?
\item Les obligations se négociaient autour de 10 le jour de l'enchère. Pourquoi le prix final
pourrait-il être inférieur~?
\item Qu'aurait-il dû se passer avec le seul règlement-livraison physique, avec davantage de protection
que d'obligations~?
\item Énoncer le \emph{résultat nommé}~: le prix final et le règlement-livraison de chaque
participant.
\item En une phrase~: qu'apporte l'enchère à un CDS~?
\end{enumerate}
\end{problem}

\section{Questions d'entretien}

\begin{interviewq}[$\star$ \iqroles{trader, chercheur}]\label{iq:m2:credit-default-swaps:1}
Qu'est-ce qu'un \omterm{def:m2:credit-default-swaps:cds}{credit default swap}~? Qui paie quoi, et quand~?
\end{interviewq}

\begin{interviewq}[$\star$ \iqroles{trader, banque}]\label{iq:m2:credit-default-swaps:2}
Pourquoi les CDS sont-ils négociés avec un coupon fixe et un \omterm{def:m2:credit-default-swaps:std}{paiement initial} (upfront)~?
\end{interviewq}

\begin{interviewq}[$\star\star$ \iqroles{chercheur}]\label{iq:m2:credit-default-swaps:3}
Dérivez la relation entre le spread d'un CDS, l'\omterm{def:m2:credit-default-swaps:hazard}{intensité de défaut} et le \omterm{def:m2:corporate-bonds:covenant}{taux de recouvrement}.
\end{interviewq}

\begin{interviewq}[$\star\star$ \iqroles{trader, chercheur}]\label{iq:m2:credit-default-swaps:4}
Comment fonctionne l'enchère de règlement-livraison d'un CDS, et pourquoi est-elle nécessaire~?
\end{interviewq}

\begin{interviewq}[$\star\star$ \iqroles{trader, chercheur}]\label{iq:m2:credit-default-swaps:5}
Qu'est-ce que la base CDS--obligation, et comment négocieriez-vous une base négative~?
\end{interviewq}

\begin{interviewq}[$\star\star\star$ \iqroles{développeur}]\label{iq:m2:credit-default-swaps:6}
Concevez un système qui valorise chaque nuit un portefeuille de 50\,000 positions sur CDS et le règle
après un \omterm{def:m2:credit-default-swaps:event}{événement de crédit}.
\end{interviewq}
